\documentclass{report}
\usepackage[T1]{fontenc} 
\usepackage[utf8]{inputenc} 
\usepackage{lmodern}
\usepackage [frenchb]{babel}
% Pour pouvoir utiliser 
\usepackage{ucs}

\usepackage{textcomp}
\usepackage{graphicx}
\usepackage{keystroke}
\usepackage{amssymb}
\usepackage{amsmath}
%\renewcommand{\thesection}{\arabic{section}} % numérotation des sections	
\usepackage[cc]{titlepic} %rajouter le logo dans la page de garde
\usepackage{url} % Pour avoir de belles url
\usepackage[a4paper]{geometry}
\usepackage[linktocpage]{hyperref}
% Pour pouvoir faire plusieurs colonnes
\usepackage {multicol}
%Pour faire plusieurs lignes
\usepackage{multirow}
\usepackage{slashbox}
% Pour mettre du code source
\usepackage {listings}
% Pour pouvoir passer en paysage
\usepackage{lscape}	

% Pour pouvoir faire plusieurs colonnes
\usepackage {multicol}

% POur crééer un index
\usepackage{makeidx}

%insertion de pdf
\usepackage{pdfpages}


%Pour taper des algorithme en pseudo code
\usepackage{algorithm,algorithmic}
%Utilisation d'acronyme
\usepackage{acronym}

\usepackage{fancyhdr}
\pagestyle{fancy}
\usepackage{lastpage} %numérotation type k/n
\renewcommand{\footrulewidth}{0.5pt} % trait horizontal
%\renewcommand{\headrulewidth}{1pt} % suppresion du trait horizontal dans l'entête
\lfoot{Projet Services}
\cfoot{\thepage \hspace{0.15cm}sur  \pageref{LastPage} } % numérotation des pages
\rfoot{A.\textsc{M} \& R.\textsc{R}}
\fancyhead[L]{\hspace{-1cm}\includegraphics[scale=0.105]{img/logouniv.png}} %logo entête !


\fancypagestyle{plain}{% 1ères pages des chapitres:
  % \fancyhead{} % supprime l’entete...
  % \renewcommand{\headrulewidth}{0pt} % ...et le filet
}




\hypersetup{
  backref=true,
  %permet d'ajouter des liens dans...
  pagebackref=true,%...les bibliographies
  hyperindex=true, %ajoute des liens dans les index.
  colorlinks=true, %colorise les liens
  breaklinks=true, %permet le retour à la ligne dans les liens trop longs
  urlcolor= blue, %couleur des hyperliens
  citecolor=	cyan,
  bookmarks=true, %créé des signets pour Acrobat
  bookmarksopen=true,
  %si les signets Acrobat sont créés,
  %les afficher complètement.
  pdftitle={Rapport Projet Services}, %informations apparaissant dans
  pdfauthor={MARGUERITE Alain},
  %dans les informations du document
  pdfsubject={Doc}
  %sous Acrobat.
}

\makeindex
%%%% debut macro pour enlever le nom chapitre %%%%
\makeatletter
\def\@makechapterhead#1{%
  \vspace*{50\p@}%
          {\parindent \z@ \raggedright \normalfont
            \interlinepenalty\@M
            \ifnum \c@secnumdepth >\m@ne
            \Huge\bfseries \thechapter\quad
            \fi
            \Huge \bfseries #1\par\nobreak
            \vskip 40\p@
}}

\def\@makeschapterhead#1{%
  \vspace*{50\p@}%
          {\parindent \z@ \raggedright
            \normalfont
            \interlinepenalty\@M
            \Huge \bfseries  #1\par\nobreak
            \vskip 40\p@
}}
\makeatother
%%%% fin macro %%%%
%Couverture 
\widowpenalty=10000
\clubpenalty=10000
\hyphenation{permet-tant}

\title{ {\huge \textsc{Distributed System}} \\Master ALMA 2\up{eme} année \\\vspace{3cm} \emph{Causal Order Message Delivery : From Logical to Physical Time}\\{\small Encadrant : M.\textsc{Achour}}\vspace{3cm}}
\titlepic{ \includegraphics[scale=0.2]{img/logouniv}}
\author{A.\textsc{Marguerite}\\ R.\textsc{Rincé}\vspace{3cm}}
\date{Université de Nantes \\ 2 rue de la Houssinière, BP92208, F-44322 Nantes cedex 03, FRANCE}
%\vspace{20cm}


\newcommand{\cf}{\emph{cf}.}


\begin{document}

\maketitle
\renewcommand{\labelitemi}{$\bullet$} 

\clearpage

\tableofcontents
\clearpage

%\chapter*{Remerciements}
\phantomsection
\addcontentsline{toc}{section}{Remerciements}
 TODO TODO TODO TODO TODO TODO TODO TODO TODO TODO TODO TODO TODO TODO TODO TODO TODO TODO TODO TODO TODO TODO

 TODO TODO TODO TODO TODO TODO TODO TODO TODO TODO TODO TODO TODO TODO TODO TODO TODO TODO TODO TODO

\chapter{Introduction}\label{chap:Intro}
\section{Présentation du sujet}
L'objectif de ce projet/TP est de réaliser un protocole de transfert de messages entre deux machines avec des listes des canaux de communication \textit{First In First Out}. La démarche proposée par l'énnonce propose les  étapes  suivantes :
  
  
  \begin{enumerate}
  \item 
    Modéliser ce protocole avec l'outil \textsc{Roméo}.
  \item
    Construire le graphe des marquages.
  \item
    Spécifier en logique temporelle: \hfill \\
    le protocole peut toujours revenir à l’état initial.
  \item 
    Expliciter les propriétés vraies.
  \item
    Programmer une implémentation : \hfill \\
    Une machine est pilotée par l'utilisateur, l'autre par le programme.
  \item
    Tests  :\hfill \\
    rendre observable les différentes transitions.
  \end{enumerate}


\begin{algorithm}
\caption{\textbf{Recherche}(nœud $N$, rectangle $R$)}
/* \textit{Recherche l'ensemble des feuilles contenues dans un rectangle requête $d$-dimensionnel à partir d'un nœud $N$, les feuilles solutions seront stockées dans $\Lambda$} */
\label{algo:recherche}
\begin{algorithmic}[1]
\IF {$N$ n'est pas une feuille}
  \FORALL{éléments $n$ de $N$ dont le MBR intersecte $R$}
    \STATE \textbf{Recherche}($n$, $R$)
  \ENDFOR
\ELSE[$N$ est une feuille] %\COMMENT{$N$ est une feuille}
  \FORALL{éléments $n$ de $N$ dont le MBR intersecte $R$}
    \STATE ajouter $n$ à $\Lambda$
  \ENDFOR
\ENDIF
\RETURN $\Lambda$
\end{algorithmic}
\end{algorithm}

\begin{algorithm}
\caption{\textbf{Ajout}(élément $E$, nœud $NR$)}
/* \textit{Ajoute un nouvel élément $E$ dans le R-tree de nœud racine $NR$} */
\label{algo:ajout}
\begin{algorithmic}[1]
\STATE On parcourt l'arbre depuis la racine $NR$ jusqu'à la feuille appropriée. À chaque niveau, on sélectionne le nœud dont le MBR nécessite le plus petit élargissement pour contenir le MBR de $E$
\STATE En cas d'égalité, on choisit le nœud ayant le MBR de plus petite \og superficie \fg{}
\IF{la feuille atteinte $L$ n'est pas pleine}
  \STATE On insère $E$ dans $L$
  \STATE On met à jour tous les MBR de $L$ jusqu'à $NR$ pour qu'ils puissent contenir le MBR de E 
\ELSE[la feuille $L$ est pleine]
  \STATE On définie $\xi$ l'ensemble constitué de tous les éléments de $L$ et du nouvel élément~$E$\\
On choisit deux éléments $e_1$ et $e_2$ tel que la distance entre $e_1$ et $e_2$ soit la plus grande de toutes les autres paires de $\xi$\\
On crée deux nœuds $l_1$ et $l_2$ contenant respectivement $e_1$ et $e_2$
  \STATE On parcourt tous les éléments restants de $\xi$ et l'on assigne chacun d'eux à $l_1$ ou $l_2$, en fonction du MBR de ces nœuds qui permet le plus petit élargissement pour contenir le MBR de l'élément considéré
  \IF{égalité de l'élargissement}
    \STATE On assigne l'élément au nœud ayant le MBR de plus petite superficie
    \IF{égalité de superficie}
      \STATE On assigne l'élément au nœud ayant le moins d'élément
    \ENDIF
  \ENDIF
  \IF[On rappelle que m est le nombre minimal d'éléments que peut contenir un nœud]{il reste $\lambda$ éléments dans $\xi$ et qu'un nœud contient $m - \lambda$ éléments}
    \STATE On assigne tous les éléments restants dans $\xi$ à ce nœud
  \ENDIF
  \STATE On met à jour le MBR des nœuds se trouvant sur le chemin de $L$ à $NR$ pour que ceux ci couvrent le MBR de $l_1$ et $l_2$
  \STATE On effectue les \og splits \fg{} éventuels des nœuds supérieurs
\ENDIF 
\end{algorithmic}
\end{algorithm}

\chapter{Logical Time}
\section{Question 1 Design a logical-time base algorithm}
Posons les variables du processus suivantes : 

\begin{itemize}
\item  \verb+ListMessages+ : représente la liste des messages qui ont été reçus mais pas délivrés.
\item \verb+m.i+ indique l'identifiant du processus à l'origine du message.
\item \verb+m.sn+ indique le prédécesseur du message envoyé par le processus émetteur.
\item \verb+m.cb+ Liste des prédécesseurs du message.
\end{itemize}

\begin{algorithm}
  \caption{\textbf{CO\_Receive}(Message m)}
  \label{algo:receive}
  \begin{algorithmic}[1]
    \IF{ $\forall$ p$\in$ m.cb $\mid$ p.sn $\leq$ del(p.i)}
    \STATE  del(m.i) ++
    \STATE  cb +=\{m.sn, m.i\}
    \STATE  CO\_Deliver(m)
    \FORALL{m\'~ $\in$ ListMessages}
    \STATE CO\_Receive2(m')
    \ENDFOR
    \ELSE
    \STATE ListMessages.add(m)
    \ENDIF
  \end{algorithmic}
\end{algorithm}


La fonction \verb+CO_Receive2(Message m)+ est la fonction suivante : 
\begin{algorithm}
  \caption{\textbf{CO\_Receive2}( Message m)}
  %/* \textit{ où pred est la liste des prédécesseurs */}
  \label{algo:receive2}
  \begin{algorithmic}[1]
    \IF{ $\forall$ p$\in$ m.cb $\mid$ p.sn $\leq$ del(p.i)}
    \STATE  del(m.i) ++
    \STATE  cb +=\{m.sn, m.i\}
    \STATE  CO\_Deliver(m)
    \STATE  ListMessage.delete(m)
    \FORALL{ m\'~ $\in$ ListMessages}
    \STATE CO\_Receive2(m\'~)
    \ENDFOR
    \ENDIF
  \end{algorithmic}
\end{algorithm}


\section{Question 2 Analyse of the algorithm}

Les valeurs que peuvent prendre \verb+cb+ sont les suivantes :

\begin{description}
\item[Minimum] \hfill \\
  C'est la valeur 1 : (\cf{} figure \ref{fig:min} : le message \verb+m1+ n'a qu'un seul prédécesseur).
\item[Maximum] \hfill \\
  C'est le nombre de processus (\cf{} figure \ref{fig:max}: le message \verb+m3+ aura trois prédécesseurs). Cette borne est possible grâce à la mise à jour de la variable.
\end{description}

\begin{figure}[htb]
  \setlength{\unitlength}{0.8mm}
  \begin{subfigure}[b]{0.5\textwidth}
    \begin{picture}(60,40)
      \put(20,20){\vector(1,0){80}}
      \put(40,20){\circle*{2}}
      \put(20,0){\vector(1,0){80}}
      \put(40,20){\vector(3,-2){29,7}}
      \put(35,22){\shortstack[l]{m1}}
    \end{picture}
    \caption{Valeur minimum}
    \label{fig:min}
  \end{subfigure}%
  ~ %add desired spacing between images, e. g. ~, \quad, \qquad etc. 
  %(or a blank line to force the subfigure onto a new line)
  \begin{subfigure}[b]{0.5\textwidth}
    \begin{picture}(60,40)
    \put(10,20){\vector(1,0){80}}
    \put(20,20){\circle*{2}}
    \put(10,22){\shortstack[l]{m1}}
    \put(20,20){\vector(3,-2){14.8}}
    \put(10,0){\vector(1,0){80}}
    \put(10,10){\vector(1,0){80}}
    \put(20,20){\vector(3,-2){14.8}}

    \put(20,20){\vector(1,-4){5}}
    \put(45,10){\circle*{2}}
    \put(37,11){\shortstack[l]{m2}}
    \put(45,10){\vector(1,2){5}}
    \put(45,10){\vector(1,-2){5}}

    \put(70,0){\circle*{2}}
    \put(62,1){\shortstack[l]{m3}}
    \end{picture}
    \caption{Valeur maximum}
    \label{fig:max}
  \end{subfigure}
  \caption{Valeurs aux bornes de cb}\label{fig:valeurs}
\end{figure}

Regardons si les propriétés citées dans le paragraphe 2.1 de l'énonce est vérifiée : 

\begin{description}
\item[Validité] Les canaux étant sûr, aucun messages de ne peut être \og perdu \fg{}. En effet il n'y a pas de parasites dans les canaux.
\item[Intégrité] Dans l'algorithme \ref{algo:receive} si un message est reçu et qu'il est délivré immédiatement alors il n'est lu qu'une fois. Au contraire s'il n'est pas délivré il est mis en attente et sera supprimé dès qu'il sera délivré (\emph{cf.} ligne de l'algorithme \ref{algo:receive2}).
\item[CO\_Delivery] Cette propriété est vérifiée grâce à la condition à la ligne de l'algorithme \ref{algo:receive}.
\item[Terminaison] On ne peut pas perdre de messages car les canaux sont sûr.
\end{description}






% LocalWords:  Logical Time logical-time algorithm ListMessages m.i
% LocalWords:  m.sn m.cb pseudo-code CO Receive p.sn del p.i cb aa of
% LocalWords:  Deliver ListMessages.add ListMessage.delete the blah
% LocalWords:  Delivery



\chapter{From Logical Time to Physical Time}

\section{Question 3 Modify your algorithm from logical time to physical time}

 blah  blah  blah  blah  blah  blah  blah  blah  blah  blah 

 blah  blah  blah  blah  blah  blah 


 blah  blah  blah  blah  blah  blah 

\chapter{Messages with Bounded Lifetime}

\section{System model : messages with a bounded lifetime}

\subsection{Question 4 Extended the previous algorithm}
\begin{algorithm}
\caption{\textbf{Send\_Time}()}
/* \textit{ Avec horloge globale */}
\label{algo:Q4a}
\begin{algorithmic}[1]
  \IF{ $\exists$ m $\in$ ListeMessages  $\mid$ \{ m.send\_time + $\Delta$ = GC.clock() \}}
      \STATE  CO\_Deliver(m)
   \FORALL{ p $\in$ m.pred }
   \IF{del(p.i) $\leq$ p.sm}
   del(p.i) $\leftarrow$ p.sn
   \ENDIF
   \ENDFOR
  \ENDIF
\end{algorithmic}
\end{algorithm}


\begin{algorithm}
\caption{\textbf{CO\_Recceive\_With\_Delay}(Message m)}
/* \textit{ Au moment où on reçoit un message */}
\label{algo:Q4b}
\begin{algorithmic}[1]
  \IF{m.send\_time + $\Delta$ $\leq$ GC.Clock}
  \STATE Delete(m)
  \ELSE[$\forall$ p$\in$ pred $\mid$ p.sn $\leq$ del(p.i)]
   \STATE  del(e.i) $\leftarrow$ e.sn
     \STATE  cb +=\{e.sn, e.i\}
     \STATE  CO\_Deliver(m)
     \FORALL{m\'~ $\in$ ListMessages}
     CO\_Receive2(m\'~})
       \ENDFOR
       \ELSE
       \STATE ListMessages.add(m)
       \ENDIF
\end{algorithmic}
\end{algorithm}


\subsection{Question 5 Analize the algorithm}
 blah  blah  blah  blah  blah  blah  blah  blah 

\begin{algorithm}
\caption{\textbf{Q5}(Message m)}
/* \textit{ Au moment où on reçoit un message */}
\label{algo:Q5}
\begin{algorithmic}[1]
\end{algorithmic}
\end{algorithm}



 blah  blah  blah  blah 



\listoffigures
% Pour finir l'interligne de 1,5
%\bibliographystyle{alpha}
%\bibliography{biblio.bib}


\end{document}
